%=====================================================================
\chapter{Joins: the Algebra and the SQL}\label{ch:joins}
%=====================================================================
\locator{4}{Part IV --- Asking Questions}

\begin{outcomes}
By the end of this chapter you will be able to:
\begin{tight}
  \item \textbf{Define} theta join, equijoin and natural join in algebra, and
        \textbf{write} each in SQL.
  \item \textbf{Distinguish} inner from outer joins, and \textbf{predict}
        which rows an inner join will drop.
  \item \textbf{Explain} why \sqlkw{NATURAL JOIN} is dangerous, using a case
        from this course's own schema.
  \item \textbf{Read} an \sqlkw{EXPLAIN QUERY PLAN} for a join and
        \textbf{say} whether it is using an index.
  \item \textbf{Account for} a hundred-fold difference between two joins that
        return identical rows.
\end{tight}
\end{outcomes}

\begin{prereq}
\chref{algebra}: a join is a selection over a product. \chref{relational} for
nulls, whose consequence for joins this chapter measures. You need
\texttt{code/m13\_joins.py} and a seeded \texttt{university.db}.
\end{prereq}

\begin{keyterms}
join $\bullet$ theta join $\bullet$ equijoin $\bullet$ natural join ($\bowtie$)
$\bullet$ inner join $\bullet$ outer join $\bullet$ left outer join $\bullet$
right outer join $\bullet$ full outer join $\bullet$ self-join $\bullet$
\sqlkw{USING} $\bullet$ \sqlkw{ON} $\bullet$ nested-loop join $\bullet$
automatic index $\bullet$ \sqlkw{EXPLAIN QUERY PLAN} $\bullet$ \sqlkw{SCAN}
$\bullet$ \sqlkw{SEARCH}
\end{keyterms}

\section{One Operator, Three Strengths}

\chref{algebra} ended on the expression
$\sigma_{\mathit{s.roll\_no}=\mathit{e.roll\_no}}(\mathit{s} \times
\mathit{e})$. It occurs so often that it has its own symbol, and three
variants that differ only in how much the condition is allowed to say.

\begin{definitionbox}[Theta join, equijoin, natural join]
A \term{theta join} $R \bowtie_{\theta} S$ is $\sigma_{\theta}(R \times S)$
for any condition $\theta$ comparing attributes of $R$ with attributes of $S$
using $=, <, \le, >, \ge$ or $\neq$.

\smallskip
An \term{equijoin} is a theta join whose condition uses only $=$. Almost every
join you will write is one.

\smallskip
A \term{natural join} $R \bowtie S$ is an equijoin on \emph{every attribute
the two relations share by name}, keeping one copy of each. It has no
condition at all --- the condition is whatever the names happen to be.
\end{definitionbox}

\begin{examplebox}[The same join, four ways in SQL]
\begin{lstlisting}[language=sql]
-- 1. the product, then the selection: chapter 12 exactly
SELECT * FROM student s, enrolment e WHERE s.roll_no = e.roll_no;

-- 2. explicit, and the form this course uses
SELECT * FROM student s JOIN enrolment e ON s.roll_no = e.roll_no;

-- 3. when the column has the same name on both sides
SELECT * FROM student JOIN enrolment USING (roll_no);

-- 4. natural: no condition written at all
SELECT * FROM student NATURAL JOIN enrolment;
\end{lstlisting}

\smallskip
The first three are one query to the optimiser --- the measurement below shows
the same plan for each. The fourth is not, on this schema, and the reason is
the most useful thing in this chapter.
\end{examplebox}

\section{What a Join Drops}

\begin{definitionbox}[Inner and outer joins]
An \term{inner join} returns only the pairs that satisfy the condition. A row
with no partner on the other side does not appear.

\smallskip
A \term{left outer join} of $R$ with $S$ returns every row of $R$; where
there is no partner in $S$, the $S$ attributes are null. A \term{right outer
join} keeps every row of $S$; a \term{full outer join} keeps every row of
both.

\smallskip
SQLite supports all three. \sqlkw{LEFT JOIN} is the one you will need, and the
question that decides it is always the same: \emph{must a row with no partner
still appear?}
\end{definitionbox}

\dsfig{%
\begin{tikzpicture}[
  cell/.style={draw=primary!35, line width=0.5pt, minimum width=1.5cm,
               minimum height=0.40cm, font=\scriptsize\ttfamily,
               text=primary!65!black, inner sep=0pt},
  hd/.style={cell, fill=primary, text=white, font=\scriptsize\bfseries},
  gone/.style={cell, fill=accent!12, text=accent},
  nul/.style={cell, text=neutral},
  ttl/.style={font=\scriptsize\bfseries, text=#1, anchor=south},
  nt/.style={font=\scriptsize\itshape, text=neutral, align=center}
]
  % section (left table)
  \node[ttl=primary] at (-5.4,0.5) {section};
  \foreach \c [count=\i] in {course, emp\_id}{
    \node[hd] at (-6.95+\i*1.5,0.2) {\c};}
  \foreach \r/\a/\b/\s in {1/CS-103/E-114/cell, 2/CS-201/E-120/cell,
                           3/CS-305/NULL/gone}{
    \node[\s] at (-5.45,-0.4*\r+0.2) {\a};
    \node[\s] at (-3.95,-0.4*\r+0.2) {\b};}

  % inner join result
  \node[ttl=greenhl!60!black] at (0.75,0.5) {\sqlkw{JOIN} instructor};
  \foreach \c [count=\i] in {course, emp\_id, name}{
    \node[hd] at (-1.5+\i*1.5,0.2) {\c};}
  \foreach \r/\a/\b/\c in {1/CS-103/E-114/Malik, 2/CS-201/E-120/Raza}{
    \node[cell] at (0.0,-0.4*\r+0.2) {\a};
    \node[cell] at (1.5,-0.4*\r+0.2) {\b};
    \node[cell] at (3.0,-0.4*\r+0.2) {\c};}
  \node[nt, text=accent, text width=4.4cm] at (1.5,-1.35)
       {CS-305 is \textbf{gone}: its emp\_id is null and null equals
        nothing};

  % left join result
  \node[ttl=secondary!70!black] at (6.95,0.5) {\sqlkw{LEFT JOIN} instructor};
  \foreach \c [count=\i] in {course, emp\_id, name}{
    \node[hd] at (4.7+\i*1.5,0.2) {\c};}
  \foreach \r/\a/\b/\c/\s in {1/CS-103/E-114/Malik/cell,
                              2/CS-201/E-120/Raza/cell,
                              3/CS-305/NULL/NULL/nul}{
    \node[\s] at (6.2,-0.4*\r+0.2) {\a};
    \node[\s] at (7.7,-0.4*\r+0.2) {\b};
    \node[\s] at (9.2,-0.4*\r+0.2) {\c};}
  \node[nt, text width=4.4cm] at (7.7,-1.55)
       {every section survives; the missing partner is filled with null};
\end{tikzpicture}}%
{The section nobody has been assigned to teach. An inner join drops it without
comment; a left join keeps it and marks what is missing. On the seeded
database this is 73 sections out of 1{,}500.}{innerleft}

\section{Measured: Plans, Nulls and the Cost of a Missing Index}

\begin{measured}[Three joins that return the same rows and one that returns none]
\textsf{\footnotesize \texttt{code/m13\_joins.py}, seeded database (4{,}000
students, 60{,}037 enrolments), SQLite 3.51.2, Intel i5-6600, Windows 11.}

\medskip
\textbf{1. Four spellings of one inner join}

\begin{tabular}{@{}L{4.6cm}R{1.9cm}R{1.7cm}L{5.6cm}@{}}
\textbf{written as} & \textbf{rows} & \textbf{ms} & \textbf{plan} \\
comma + \sqlkw{WHERE}  & 60{,}037 & 29.1 & \texttt{SCAN e} + \texttt{SEARCH s} \\
\sqlkw{JOIN} \dots\ \sqlkw{ON} & 60{,}037 & 30.4 & identical \\
\sqlkw{JOIN} \dots\ \sqlkw{USING} & 60{,}037 & 30.1 & identical \\
\sqlkw{NATURAL JOIN} & \textbf{0} & 43.7 & \\
\end{tabular}

\smallskip
The two tables share \textbf{two} attribute names, not one:
\texttt{roll\_no}, and \texttt{semester} --- which is an \texttt{INTEGER} year
of study (e.g.\ \texttt{8}) in \textit{student} and a \texttt{TEXT} term label
(e.g.\ \texttt{'FALL2025'}) in \textit{enrolment}. \sqlkw{NATURAL JOIN} matched
on both, demanded a year equal a term label, and returned nothing. \textbf{No
error, on the schema exactly as \chref{mapping} built it.}

\medskip
\textbf{2. Inner against left join}

\begin{tabular}{@{}L{9.4cm}R{2.4cm}@{}}
sections in the database & 1{,}500 \\
sections with no instructor yet (\texttt{emp\_id} null) & 73 \\
timetable as an inner join & 1{,}427 \\
timetable as a left join & 1{,}500 \\
\end{tabular}

\smallskip
The 73 missing from the inner join are exactly the sections nobody has been
assigned to --- the rows a timetable most needs to show.

\medskip
\textbf{3. Students who share a name: a join on an unindexed column}

\begin{tabular}{@{}R{1.6cm}R{2.0cm}R{1.9cm}R{2.4cm}R{2.0cm}R{1.8cm}@{}}
\textbf{students} & \textbf{pairs found} & \textbf{as written}
& \textbf{automatic index off} & \textbf{real index} & \textbf{slowdown} \\
1{,}000 & 1{,}785  & 3.2\,ms  & 235.1\,ms  & 2.9\,ms  & 82$\times$ \\
2{,}000 & 7{,}109  & 8.6\,ms  & 1{,}088.1\,ms & 11.1\,ms & 98$\times$ \\
4{,}000 & 28{,}517 & 25.1\,ms & 4{,}170.1\,ms & 38.7\,ms & \textbf{108$\times$} \\
\end{tabular}

\smallskip
Without an index the join is a \term{nested-loop join}: for each row of
one side, every row of the other is examined --- two loops, one inside the
other, which is where the $n^2$ comes from.

\smallskip
\textbf{Compare across a row, not down a column.} The pairs found grow about
four-fold per doubling for \emph{every} method, because names come from a
fixed pool and the number of real matches is itself quadratic --- so growth
ratios here mix up how much the database \emph{looked at} with how much it
\emph{returned}. Across a row the output is identical, and the difference is
pure search: without an index, all $4000 \times 3999 / 2 \approx 8$ million
pairs are examined to find 28 thousand. The slowdown itself rises with size,
82 to 108, which is the real signature of a quadratic method.

\medskip
\textbf{4. Why ``as written'' is fast anyway}

\begin{tabular}{@{}L{4.8cm}L{9.2cm}@{}}
automatic index on (default) & \texttt{SEARCH b USING AUTOMATIC COVERING INDEX (name=?)} \\
automatic index off & \texttt{SEARCH b USING INDEX \dots (roll\_no>?)} \\
\texttt{CREATE INDEX} on name & \texttt{SEARCH b USING INDEX idx\_name (name=?)} \\
\end{tabular}

\smallskip
SQLite noticed the missing index, \textbf{built one for this single query,
used it, and threw it away}. You pay to build an index on every run and never
get to keep it.
\end{measured}

\begin{alertbox}[What the measurement means for the queries you write]
\begin{tight}
  \item \textbf{Never use \sqlkw{NATURAL JOIN}.} Its condition is decided by
        column names, and column names change for reasons that have nothing to
        do with the join. This course's own schema has the trap in it, and it
        was not put there on purpose --- a year of study and a term are both
        reasonably called ``semester''. Write \sqlkw{ON} or \sqlkw{USING} and
        say what you mean.
  \item \textbf{Ask, for every join, whether an unmatched row must still
        appear.} If yes, it is a \sqlkw{LEFT JOIN}. The inner join's silence
        about the 73 is not a bug in SQL; it is the definition, applied to a
        null foreign key that \chref{cardinality} deliberately allowed.
  \item \textbf{Read the plan.} \sqlkw{AUTOMATIC INDEX} is the optimiser
        telling you an index is missing, in plain words. It is the cheapest
        performance diagnosis there is, and \chref{indexes} is about acting on
        it.
\end{tight}
\end{alertbox}

\section{Joining a Table to Itself}

\begin{examplebox}[The self-join, and the condition that halves it]
``Pairs of students who share a name'' joins \textit{student} to itself, which
needs two aliases --- \chref{algebra}'s rename:

\begin{lstlisting}[language=sql]
SELECT a.roll_no, b.roll_no, a.name
  FROM student a
  JOIN student b ON a.name = b.name
                AND a.roll_no < b.roll_no;
\end{lstlisting}

\smallskip
\textbf{The second condition does two jobs.} Without it every student matches
themselves --- Ayesha shares a name with Ayesha --- and every real pair appears
twice, once each way round. \texttt{a.roll\_no < b.roll\_no} removes both: no
row is paired with itself, and each pair appears in exactly one order. Use
$<$, never $\neq$, which would keep both orders.

\smallskip
The prerequisite query in \chref{algebra}'s laboratory was a self-join too, of
\textit{course}, through \textit{prerequisite}. Every recursive relationship in
an ER diagram becomes a self-join in SQL.
\end{examplebox}

\begin{pitfall}
\begin{tight}
  \item \textbf{\sqlkw{NATURAL JOIN}.} It joined on two columns, not one, and
        returned nothing. Measured.
  \item \textbf{An inner join where a left join was meant.} The unmatched rows
        vanish silently --- 73 sections from a timetable.
  \item \textbf{A condition on the right table in \sqlkw{WHERE} after a
        \sqlkw{LEFT JOIN}.} \texttt{WHERE i.name = 'Malik'} discards the
        null-extended rows and turns it back into an inner join. Put such
        conditions in \sqlkw{ON}.
  \item \textbf{Forgetting the join condition.} You get the product: 240
        million rows on the seeded data.
  \item \textbf{A self-join with $\neq$.} Every pair twice. Use $<$.
  \item \textbf{Ignoring \sqlkw{AUTOMATIC INDEX} in a plan.} It is a missing
        index announcing itself.
\end{tight}
\end{pitfall}

\begin{lab}[Lose 73 sections, find them, and make a join a hundred times faster]
\begin{lstlisting}[language=bash]
# --- 1. The seeded database ----------------------------------------
python code/schema/seed.py university.db
sqlite3 university.db
\end{lstlisting}
\begin{lstlisting}[language=sql]
-- --- 2. The timetable, as an inner join -------------------------
.headers on
.mode box
SELECT COUNT(*) FROM section s JOIN instructor i ON s.emp_id = i.emp_id;
SELECT COUNT(*) FROM section;
-- The two numbers differ. By how much, and which sections are they?

-- --- 3. Find them -----------------------------------------------
SELECT s.course_code, s.section_no, s.semester
  FROM section s LEFT JOIN instructor i ON s.emp_id = i.emp_id
 WHERE i.emp_id IS NULL
 LIMIT 10;
-- This is the anti-join idiom: LEFT JOIN, then keep the rows that
-- found no partner. Learn it; Chapter 15 measures it.

-- --- 4. Break the left join by mistake --------------------------
SELECT COUNT(*) FROM section s LEFT JOIN instructor i
    ON s.emp_id = i.emp_id
 WHERE i.designation = 'Professor';
SELECT COUNT(*) FROM section s LEFT JOIN instructor i
    ON s.emp_id = i.emp_id AND i.designation = 'Professor';
-- Same condition, two places, two very different counts. Explain.

-- --- 5. The natural join, on the schema you built ---------------
SELECT COUNT(*) FROM student NATURAL JOIN enrolment;
-- Zero. Find out why:
PRAGMA table_info(student);
PRAGMA table_info(enrolment);
-- Name the second shared column, and its two meanings.

-- --- 6. Read a plan ---------------------------------------------
EXPLAIN QUERY PLAN
SELECT COUNT(*) FROM student a JOIN student b
    ON a.name = b.name AND a.roll_no < b.roll_no;
-- Find the word AUTOMATIC. Say what SQLite just did.

-- --- 7. Give it the index it keeps building ---------------------
CREATE INDEX idx_student_name ON student(name);
EXPLAIN QUERY PLAN
SELECT COUNT(*) FROM student a JOIN student b
    ON a.name = b.name AND a.roll_no < b.roll_no;
\end{lstlisting}
\begin{lstlisting}[language=bash]
# --- 8. The measured version ----------------------------------------
python code/m13_joins.py
\end{lstlisting}

\smallskip
\textbf{What to notice.} Step 4 is the one that catches people in examinations
and in production alike. A condition in \sqlkw{WHERE} is applied \emph{after}
the outer join has produced its null-extended rows, and since a null
\texttt{designation} is not \texttt{'Professor'}, they are all discarded ---
the outer join has silently become an inner one.
\end{lab}

\begin{checkpoint}
\begin{tightnum}
  \item \sqlkw{NATURAL JOIN} of \textit{student} and \textit{enrolment}
        returns 0 rows on this course's schema. Explain exactly why, say what
        you would change in the query, and say what you would change in the
        schema if you could change only one column name.
  \item An inner join of \textit{section} to \textit{instructor} returns 1{,}427
        rows and the table has 1{,}500. Identify the missing 73 by a property
        they share, write the query that lists them, and say which design
        decision in \chref{cardinality} made them possible.
  \item Two joins return the same 28{,}517 rows, one in 38.7\,ms and one in
        4{,}170\,ms. Explain the difference using the number of pairs each one
        examined, and say why comparing growth ratios down the columns would
        have given a misleading answer.
\end{tightnum}
\end{checkpoint}

\begin{chaptersummary}
\begin{tight}
  \item A \textbf{theta join} is $\sigma_\theta(R \times S)$; an
        \textbf{equijoin} uses only $=$; a \textbf{natural join} equijoins on
        \emph{every} shared attribute name.
  \item Comma-plus-\sqlkw{WHERE}, \sqlkw{JOIN \dots ON} and \sqlkw{JOIN \dots
        USING} produce the same plan. Use \sqlkw{ON}.
  \item \textbf{Measured: \sqlkw{NATURAL JOIN} returned 0 rows} on this
        course's schema, because \textit{student} and \textit{enrolment} both
        have a \texttt{semester} --- a year of study in one, a term label in
        the other. Never use it.
  \item An \textbf{inner join} drops rows with no partner; a \textbf{left
        join} keeps them with nulls. \textbf{Measured: 73 of 1{,}500 sections}
        vanished from the inner-join timetable --- the ones with no instructor.
  \item A condition on the right-hand table belongs in \sqlkw{ON}. In
        \sqlkw{WHERE} it turns a left join back into an inner one.
  \item A \textbf{self-join} needs two aliases, and $<$ rather than $\neq$ to
        avoid both self-pairs and duplicate pairs.
  \item \textbf{Measured: a join on an unindexed column was 108$\times$ slower}
        at 4{,}000 students than the same join with an index --- 8 million
        pairs examined against 28 thousand found.
  \item Compare methods \textbf{at the same size}. Here the output itself grew
        quadratically, so growth ratios mixed search cost with output size.
  \item SQLite builds an \textbf{automatic index} for a single query when one
        is missing, and discards it. \sqlkw{AUTOMATIC} in a plan is a missing
        index announcing itself.
\end{tight}
\end{chaptersummary}

\begin{reviewq}
  \item \textbf{(MCQ)} A natural join of $R$ and $S$ is an equijoin on:
        (a)~their primary keys; (b)~their foreign keys; (c)~every attribute
        they share by name; (d)~the first column of each.
  \item \textbf{(MCQ)} A section whose \texttt{emp\_id} is null, joined to
        \textit{instructor} with an inner join, appears: (a)~once, with nulls;
        (b)~not at all; (c)~once per instructor; (d)~with an error.
  \item \textbf{(Short)} Why does \texttt{a.roll\_no < b.roll\_no} in a
        self-join give each pair exactly once, while $\neq$ gives it twice?
  \item \textbf{(Short)} What does \sqlkw{AUTOMATIC COVERING INDEX} in a query
        plan tell you, and what should you do about it?
  \item \textbf{(Short)} Give a query in the university schema that needs a
        \sqlkw{LEFT JOIN}, and say what an inner join would silently omit.
  \item \textbf{(Applied)} Write the query that lists every course with its
        number of sections in FALL2026, including courses with none. (On the
        seeded database every course has one, so first insert a new course
        with no sections, or there is nothing to omit.) Show the inner-join
        version that gets it wrong, say which courses it omits, and then give
        the correct version and explain where its condition on semester must
        go and why.
  \item \textbf{(Applied)} The measured slowdown was 82$\times$, 98$\times$ and
        108$\times$ at 1{,}000, 2{,}000 and 4{,}000 students. Predict it at
        8{,}000, show your reasoning from the number of pairs examined and
        found, and explain why the slowdown is rising rather than constant.
\end{reviewq}
